\xhdr{Our contributions} 
In detail, our contributions are as follows.

$\bullet$ We initiate the study of low-distortion mechanisms for committee selection that incorporate \emph{fairness} constraints under the egalitarian (min-max) objective for the social cost, by modeling the task as the ordinal fair $k$-center problem with limited access to cardinal information. 

$\bullet$ We present a $5$-distortion algorithm with only $\bigO(k \log ^2 k)$ cardinal queries. 
Additionally, we present a $3$-distortion algorithm that requires $2k^2$ queries. 

\subsection{Our techniques}
Our methods build upon techniques from both the ordinal setting with limited query access and the clustering domain.
The ordinal $k$-center under limited cardinal information was studied by
\citet{burkhardt2024low}, where they presented two algorithmic results: a $2$-distortion algorithm with $\frac{k^2-k}{2}$ queries and a $4$-distortion algorithm with only $2k$ queries. Both these algorithms are based on adaptions of~\citet{gonzalez1985clustering}'s algorithm.
On the other hand, in the clustering domain,
a common framework that has been successful in solving the fair $k$-center problem (with complete cardinal information), is that of maximal matching~\cite{jones2020fair,gadekar2025fair}.
This approach, on a high level,  first solves the standard $k$-center problem to obtain a good quality set of ``fairness-oblivious'' centers $T$ of size $k$ with good approximation guarantee, and
projects it onto a set of fair centers via maximal matching, inuring only a small loss in the quality.
In our work, we take the set $T$ to be the output of the algorithm of~\cite{burkhardt2024low}, using either the $2$-distortion or $4$-distortion variant. We first describe the underlying idea for obtaining a fair set of centers from $T$ via the maximum matching framework of~\cite{jones2020fair,gadekar2025fair}. We then discuss the main challenges in adapting this framework to the ordinal setting, and finally present our novel techniques for overcoming these challenges.
For clarity of presentation, we restrict our attention to the case with $k$ groups, $\mathbb{G}=\{G_1,\dots,G_k\}$, each having a unit requirement, and the optimal cost is $1$. 
Let $d$ be the distance function of the underlying metric space and $n$ be the number of points (or agents) in the input instance.

The key idea of the matching-based approach consists of two components. 
The first component is to use the properties of the (ordered) set $T=(t_1,\dots,t_k)$ adapted from the approaches of~\citet{gonzalez1985clustering}.  
In particular, such set $T$ has a \emph{critical} index $\ell^* \in [k]$, which is the largest index $\ell \in [k]$ such that the points $T_\ell := (t_1,\dots,t_\ell)$ of $T$ are from different clusters of an (fixed) optimal solution, and the cost of $T_\ell$ is within $\alpha$ factor of the optimal cost, for some constant $\alpha>1$. 
The second component is to project the ``fairness-oblivious'' centers $T_{\ell}$ onto a set of fair centers via maximal matching. 
In more detail, for $\ell \in [k]$, consider the bipartite graph $H_\ell$ between the centers of $T_{\ell}$ (left partition) and the demographic groups $\GroupSet$ (right partition), with edges between $t \in T_{\ell}$ and group $G \in \GroupSet$ if $d(t,G) \le \lambda$, for some non-negative real $\lambda$.
Let $\lambda_{\ell}$ be the minimum $\lambda$ such that there exists a left-perfect matching on $H$, noting that $\lambda_\ell\le 1$.
This implies that, if we find a left-perfect matching on $H_\ell$, then we can obtain a set of fair centers of size $k$ from $T_{\ell}$.
Furthermore, the cost of such a solution, denoted as $Sol(\ell)$, is at most $\cost(T_{\ell}) + \lambda_{\ell}$, by triangle inequality.
Moreover, when $\ell = \ell^*$, $H_{\ell^*}$ has a perfect matching on the left-partition ($T_{\ell^*}$), and hence $T_{\ell^*}$ can be extended to a fair solution with cost $\alpha+1$.



Notice that, $\lambda_{\ell^*}$ belongs to one of the $nk$ distances, and $\ell^{*}$ can be guessed over the range $[k]$.
This immediately implies an algorithm to obtain distortion $3$ (and distortion $5$ resp.), by combining the $2$-distortion (and $4$-distortion resp.) algorithm of \cite{burkhardt2024low} for $T$ with the matching frameworks of~\cite{jones2020fair,gadekar2025fair}, using additional $nk(k\log (nk))$ distance queries, where the $k \log (nk)$ factor is for constructing the (edges of the) bipartite graph. 
This naturally leads to the following question: \emph{can we achieve the same distortion guarantees with smaller number of queries, preferably near-linear in $k$, to the cardinal information?}

\vspace{1em}



For each $\ell \in [k]$, one possible idea is to employ binary search on $nk$ distances to find $\lambda_{\ell}$. However, the number of queries is still linear in $nk$, due to the need to sort these distances for the binary search.
To reduce the number of queries, our first idea is to shrink the search space of candidate distances from $nk$ to $k^2$. For each center $t \in T$ and each group $G$ (there are at most $k$ groups), we identify the nearest point in $G$ to $t$ using only the ordinal ranking of $t$. 
This results in at most $k^2$ distances, and hence requires $k^2$ many queries to construct the bipartite graphs that facilitate us find the minimum $\lambda_{\ell}$ for any $\ell \in [k]$.
This yields a $3$-distortion algorithm using $\bigO(k^2)$ queries.
We give a detailed analysis in Section~\ref{sec:conventional-algorithms}. 



A natural way to reduce the distance queries is to conduct a binary search on $\ell$ over function $Sol(\ell)$. However, this not possible:  as $cost(T_{\ell})$ is non-increasing and $\lambda_{\ell}$ is non-decreasing, $Sol(\ell)$ is not necessarily monotone. 
Our first key contribution is introducing the predicate $$P(\ell) \equiv (4 \lambda_{\ell} \leq cost(T_{\ell}));$$
which \emph{is monotone} and facilitates binary search. We find the largest $\ell$, denoted $L$, such that $P(L)$ is true but $P(L+1)$ is false. 
Using properties we derive in Lemma~\ref{lemma:monotonicity-of-lambda-and-cost-init-sol-set}, we prove that $\min\{Sol(L), Sol(L+1)\} \leq 5 cost(S^*)$ where $S^*$ is the optimal solution. This establishes the correctness of our binary-search procedure. 

Next, we handle the distance query complexities of evaluating $P(\ell)$ and $\lambda_{\ell}$, which are completely new problems. 

For the query complexity of computing $\lambda_{\ell}$ (detailed in Section~\ref{section:finding-minimum-lambda}), at a high level, we reduce the search space for $\lambda_{\ell}$ (of size $k^2$) by a factor of $\frac{3}{4}$ at each iteration. To accomplish this, at each round, we identify a \emph{pivot point} in the search space such that at least $\frac{1}{4}$ of the candidate distances are no smaller than the pivot and at least $\frac{1}{4}$ are no greater than the pivot such that at least one part does not contain $\lambda_\ell$. Finally, using a median-of-medians-style subroutine (Algorithm~\ref{alg:finding-lambda-ell} and Algorithm~\ref{alg:median-median}), we find such a pivot; this yields query complexity $O(\ell \log^{2} k)$.
Notice that, for the whole algorithm, we only need to compute $\lambda_{\ell}$ for a fixed $\ell$ twice. Therefore, in the worst case scenario when $\ell = k$, the total query complexity of computing $\lambda_{\ell}$ is $O(k \log^2 k)$.

For the query complexity of evaluating $P(\ell)$, a natural approach is to compute  $\lambda_{\ell}$ and compare it with $T_{\ell}$; however, this approach needs to compute $\lambda_{\ell}$ $\log(k)$ times and leads to  $O(\ell \log^3(k))$ queries. To avoid repetitively computing $\lambda_{\ell}$, we apply the equivalent formulation of $P(\ell)$ as 
$$P(\ell)\equiv (\lambda_{\ell} \leq \frac{cost(T_{\ell})}{4}),$$
which holds precisely when there exists a left-perfect matching when setting $\lambda = \frac{cost(T_{\ell})}{4}$. Using this design, we obtain an overall query complexity of $O(k \log^2 k)$. 




